\section{One orbit point per moving row}
\label{sec:orbit-sparsity}

The complement \(\sigma\) in \eqref{eq:opened-factorization} is not a
free affine slope.  It remembers the physical target:
\[
 m_\alpha j+h_0=u\sigma.
\]
This provenance converts conductor degeneracy into orbit sparsity.

\begin{theorem}[Orbit sparsity on the constant branch]
\label{thm:orbit-sparsity}
Fix a moving physical row \(\alpha\).  Among all orbit points
\(j\in\cJ_X\), at most one can support a conductor-one key
obtained by opening the ultra divisor on that row.
\end{theorem}

\begin{proof}
By \cref{thm:constant-dichotomy}, a conductor-one key satisfies
\(q_X\mid\sigma\).  The factorization
\eqref{eq:opened-factorization} then gives
\begin{equation}
 m_\alpha j+h_0\equiv0\pmod{q_X}.
\label{eq:orbit-congruence}
\end{equation}
Since \(q_X\nmid m_\alpha\), this congruence has exactly one residue
class
\[
 j\equiv-h_0m_\alpha^{-1}\pmod{q_X}.
\]
By \cref{lem:padding-consequences}, \(\cJ_X\) contains at most one
integer in that class.
\end{proof}

\begin{corollary}[Both polarizations]
\label{cor:both-polarizations-sparse}
For the left-ultra polarization, each left moving row admits at most
one constant-sector orbit point while the right row remains arbitrary.
For the right-ultra polarization, the symmetric statement holds with
the row labels exchanged.
\end{corollary}

\begin{remark}[Divisor choices do not restore the \(J_X\) factor]
\label{rem:divisor-choices}
At the unique possible orbit point, several divisors \(u\) can make
\(\sigma=(m_\alpha j+h_0)/u\) divisible by \(q_X\).  Their total
absolute coefficient is \(X^{o(1)}\) by
\eqref{eq:ultra-envelope}.  Thus divisor multiplicity is a soft
factor; it cannot restore \(J_X\) distinct orbit points.
\end{remark}

\begin{proposition}[Sharpness of the one-point statement]
\label{prop:one-point-sharp}
The conclusion of \cref{thm:orbit-sparsity} cannot be improved to
zero using only interval length and primitivity.  For any prime
\(q\nmid m\) and any fixed \(h_0\), an interval containing
\(-h_0m^{-1}\pmod q\) has one solution of
\(q\mid mj+h_0\).
\end{proposition}

\begin{proof}
Choose the canonical representative of the unique residue class and
an interval containing it.  No M\"obius sign or target mask has been
used, so nothing in the stated hypotheses forbids that point.
\end{proof}
